\documentclass[11pt]{article}

\usepackage[margin=1in]{geometry}
\usepackage{amsmath,amssymb,amsthm}
\usepackage{algorithm}
\usepackage{algpseudocode}
\usepackage{booktabs}
\usepackage{hyperref}

\newcommand{\R}{\mathbb{R}}
\newcommand{\NN}{\mathcal{N}}
\newcommand{\CC}{\mathcal{C}}
\newcommand{\FF}{\mathcal{F}}

\title{Classical Ruge--St\"uben Algebraic Multigrid:\\
  Strength of Connection and Standard Interpolation}
\author{}
\date{}

\begin{document}
\maketitle

\section{Introduction}

Algebraic multigrid (AMG) methods solve sparse linear systems $Ax = b$
by constructing a hierarchy of progressively coarser problems without
requiring geometric grid information.  The classical AMG method of Ruge
and St\"uben~\cite{RS1987} achieves this through three key components:
\begin{enumerate}
  \item \textbf{Strength of connection} --- identifying which matrix
    entries represent important couplings between unknowns.
  \item \textbf{Coarse/fine (C/F) splitting} --- partitioning unknowns
    into coarse points (retained on the next level) and fine points
    (eliminated by interpolation).
  \item \textbf{Standard interpolation} --- constructing the prolongation
    operator $P$ that maps coarse-grid corrections to the fine grid.
\end{enumerate}
Given $P$, the restriction operator is $R = P^T$ and the coarse-grid
operator is formed by the Galerkin product $A_c = P^T A P$.

Throughout, we assume $A \in \R^{n \times n}$ is an M-matrix: $a_{ii} > 0$,
$a_{ij} \le 0$ for $j \ne i$.  We denote the set of neighbors of point~$i$
(off-diagonal entries in row~$i$) by $\NN_i = \{j \ne i : a_{ij} \ne 0\}$.

%======================================================================
\section{Strength of Connection}
\label{sec:strength}
%======================================================================

The fundamental observation underlying AMG is that relaxation methods
(e.g., Gauss--Seidel) effectively reduce oscillatory error components
but leave \emph{algebraically smooth} error --- error that satisfies
$Ae \approx 0$ pointwise --- largely unchanged~\cite{RS1987,Stuben2001}.
For the $i$-th equation,
\[
  a_{ii}\,e_i + \sum_{j \in \NN_i} a_{ij}\,e_j \approx 0
\]
implies that $e_i$ depends most strongly on neighbors $j$ for which
$|a_{ij}|$ is large relative to the other couplings in the row.

\begin{definition}[Strong connection]
  Given a threshold parameter $\theta \in (0,1]$, point~$j$ is a
  \textbf{strong connection} of point~$i$ if
  \begin{equation}
    \label{eq:strength}
    -a_{ij} \;\ge\; \theta \cdot \max_{k \ne i}\bigl(-a_{ik}\bigr).
  \end{equation}
  The set of strong connections of~$i$ is denoted $S_i$.  We also
  define the \textbf{strong transpose}:
  \[
    S_i^T = \{j : i \in S_j\}
  \]
  which is the set of points that~$i$ strongly influences.
\end{definition}

Typical values are $\theta = 0.25$ for two-dimensional problems and
$\theta = 0.5$ for three-dimensional problems~\cite{Stuben2001}.  For
an M-matrix, $-a_{ij} \ge 0$ for all $j \ne i$, so the
criterion~\eqref{eq:strength} selects off-diagonal entries whose
magnitude is at least a fraction~$\theta$ of the row maximum.

The strength matrix $S$ is stored in compressed sparse row (CSR) format
containing only the strong off-diagonal entries.  Construction requires
three passes over~$A$: (i)~compute the row-wise maximum, (ii)~count
strong entries per row, (iii)~fill the CSR arrays.  The total cost is
$O(\mathrm{nnz}(A))$.

%======================================================================
\section{C/F Splitting}
\label{sec:coarsening}
%======================================================================

The coarsening step partitions $\{1, \ldots, n\}$ into a set of
\textbf{coarse points}~$\CC$ and \textbf{fine points}~$\FF$.  The
splitting should satisfy two heuristic principles~\cite{RS1987}:

\begin{description}
  \item[(H1)] For each F-point~$i$, every point $j \in S_i$ is either
    a C-point or is strongly connected to a point in $\CC_i = S_i \cap \CC$.
  \item[(H2)] The C-points form a maximal independent set with respect
    to the strong connections (no two C-points are strongly connected).
\end{description}

Heuristic~(H1) ensures that the interpolation formula (Section~\ref{sec:interp})
can account for all important couplings of each F-point.
Heuristic~(H2) keeps the coarse grid as small as possible.
In practice~(H2) is treated as a soft guideline; the coarsening ratio
$n / |\CC|$ is typically between 2 and 3 for structured grids.

\subsection{First Pass: Greedy Independent Set}

The algorithm uses the \emph{influence measure}
\begin{equation}
  \label{eq:lambda}
  \lambda_i = \bigl|\{j \in S_i^T : j \text{ is undecided}\}\bigr|,
\end{equation}
which counts how many undecided points are strongly influenced by~$i$.
Computing $\lambda$ requires the transpose $S^T$.

\begin{algorithm}[H]
\caption{Classical RS Coarsening --- First Pass}
\label{alg:coarsen}
\begin{algorithmic}[1]
  \State Initialize $\lambda_i \gets |S_i^T|$ and $\mathrm{cf}[i]
    \gets \textsc{Undecided}$ for all~$i$
  \While{undecided points remain}
    \State Select $i^* = \arg\max_i \{\lambda_i : \mathrm{cf}[i] =
      \textsc{Undecided}\}$
    \State $\mathrm{cf}[i^*] \gets C$ \hfill
      \Comment{$i^*$ becomes a coarse point}
    \For{each $j \in S_{i^*}^T$ with $\mathrm{cf}[j] = \textsc{Undecided}$}
      \State $\mathrm{cf}[j] \gets F$ \hfill
        \Comment{influenced points become fine}
      \For{each $k \in S_j$ with $\mathrm{cf}[k] = \textsc{Undecided}$,
        $k \ne i^*$}
        \State $\lambda_k \gets \lambda_k + 1$ \hfill
          \Comment{$k$ gains value: $j$ is now F and needs a C-neighbor}
      \EndFor
    \EndFor
  \EndWhile
\end{algorithmic}
\end{algorithm}

The increment on line~8 is the key mechanism: when $j$ becomes an
F-point, it will need C-points nearby for interpolation.  The undecided
points that $j$ depends on ($k \in S_j$) become more attractive as
future C-points, so their $\lambda$ increases.  This biases subsequent
selections toward fulfilling heuristic~(H1).

\subsection{Second Pass}

After the first pass, rare configurations may leave an F-point with no
strong C-neighbor.  The second pass promotes any such F-point to~C:
\[
  \text{If } i \in \FF \text{ and } S_i \cap \CC = \emptyset,
  \quad\text{then } i \to \CC.
\]

\subsection{Efficient Implementation via Bucket Structure}

A naive implementation scans all $n$ points to find $\arg\max \lambda$
at each step, giving $O(n^2)$ total cost.  This is replaced by a
\emph{bucket structure}~\cite{Stuben2001,HensonYang2002} that maintains
the priority queue in $O(1)$ per operation.

\paragraph{Data structure.}
Maintain an array of doubly-linked lists indexed by $\lambda$ value:
\begin{itemize}
  \item $\texttt{bucket\_head}[k]$: head of the list of points with
    $\lambda = k$, for $k = 0, 1, \ldots, B-1$.
  \item $\texttt{next}[i]$, $\texttt{prev}[i]$: forward and backward
    pointers for point~$i$ within its bucket list ($-1$ as sentinel).
  \item $\texttt{top\_bucket}$: index of the highest non-empty bucket.
\end{itemize}
The number of buckets $B$ is bounded by the initial maximum $\lambda$
plus the maximum row length of~$S$ (since $\lambda$ values only increase
via the increment on line~8 of Algorithm~\ref{alg:coarsen}, and each
point's $\lambda$ is incremented at most once per F-point creation in
its neighborhood).

\paragraph{Operations.}
All three priority-queue operations run in $O(1)$ amortized time:
\begin{enumerate}
  \item \textsc{FindMax}: Return $\texttt{bucket\_head}[\texttt{top\_bucket}]$.
    If the bucket is empty, decrement $\texttt{top\_bucket}$ until a
    non-empty bucket is found.  Over the entire algorithm, $\texttt{top\_bucket}$
    decreases at most $B$ times total, so the amortized cost is $O(1)$.
  \item \textsc{Remove}$(i)$: Unlink~$i$ from its doubly-linked list.
    Constant time.
  \item \textsc{Move}$(i, \lambda_{\mathrm{old}} \to \lambda_{\mathrm{new}})$:
    Remove from bucket $\lambda_{\mathrm{old}}$, insert at the head of
    bucket $\lambda_{\mathrm{new}}$, update $\texttt{top\_bucket}$ if
    $\lambda_{\mathrm{new}} > \texttt{top\_bucket}$.  Constant time.
\end{enumerate}

When a point is selected as C or marked as F, it is removed from its
bucket and not reinserted.  When $\lambda_k$ is incremented (line~8),
point~$k$ is moved from bucket $\lambda_k$ to bucket $\lambda_k + 1$.

\paragraph{Total complexity.}
Each point is inserted, removed, and possibly moved a bounded number of
times.  Each edge of~$S$ is traversed at most twice (once via $S^T$ when
the C-point is selected, once via~$S$ when the F-point increments its
neighbors).  The total work is therefore
$O(n + \mathrm{nnz}(S))$.

%======================================================================
\section{Standard (Direct) Interpolation}
\label{sec:interp}
%======================================================================

The prolongation operator $P \in \R^{n \times n_c}$ is defined row by
row.  For a C-point~$i$ with coarse index~$\hat\imath$, row~$i$ of~$P$
has a single entry:
\[
  P_{i,\hat\imath} = 1 \qquad\text{(identity injection)}.
\]
For an F-point~$i$, the weights are derived from the algebraic
smoothness assumption.

\subsection{Neighbor Classification}

For F-point~$i$, partition its neighbors into three sets:
\begin{center}
\begin{tabular}{lll}
  \toprule
  Set & Definition & Role \\
  \midrule
  $\CC_i$ & $\{j \in S_i : j \in \CC\}$ &
    Strong C-neighbors (interpolation sources) \\
  $D_i^s$ & $\{k \in S_i : k \in \FF\}$ &
    Strong F-neighbors (redistributed to $\CC_i$) \\
  $D_i^w$ & $\NN_i \setminus S_i$ &
    Weak connections (lumped into diagonal) \\
  \bottomrule
\end{tabular}
\end{center}

\subsection{Derivation}

Start from the smoothness condition for row~$i$:
\begin{equation}
  \label{eq:smooth}
  a_{ii}\,e_i + \sum_{j \in \CC_i} a_{ij}\,e_j
    + \sum_{k \in D_i^s} a_{ik}\,e_k
    + \sum_{m \in D_i^w} a_{im}\,e_m \;\approx\; 0.
\end{equation}

\paragraph{Step 1: Lump weak connections.}
For $m \in D_i^w$, the connection to~$i$ is weak, so the error varies
little: $e_m \approx e_i$.  Substituting and collecting:
\begin{equation}
  \label{eq:lump}
  \hat{a}_{ii}\,e_i + \sum_{j \in \CC_i} a_{ij}\,e_j
    + \sum_{k \in D_i^s} a_{ik}\,e_k \;\approx\; 0,
\end{equation}
where the \textbf{modified diagonal} is
\begin{equation}
  \label{eq:diaghat}
  \hat{a}_{ii} = a_{ii} + \sum_{m \in D_i^w} a_{im}.
\end{equation}

\paragraph{Step 2: Redistribute strong F-neighbors.}
For $k \in D_i^s$, approximate $e_k$ by interpolation from the
C-points in~$\CC_i$:
\begin{equation}
  \label{eq:redistrib}
  e_k \;\approx\; \sum_{j \in \CC_i}
    \frac{a_{kj}}{\displaystyle\sum_{\ell \in \CC_i} a_{k\ell}}\;e_j.
\end{equation}
This distributes the contribution of F-neighbor~$k$ to the C-points
in~$\CC_i$, weighted proportionally to the strength of connection from
$k$ to each C-point.  If $k$ has no connection to any point in~$\CC_i$
(the denominator in~\eqref{eq:redistrib} vanishes), then $a_{ik}$ is
added to $\hat{a}_{ii}$ instead.

\paragraph{Step 3: Solve for interpolation weights.}
Substituting~\eqref{eq:redistrib} into~\eqref{eq:lump}:
\[
  \hat{a}_{ii}\,e_i + \sum_{j \in \CC_i}
    \underbrace{\left(a_{ij} + \sum_{k \in D_i^s}
      \frac{a_{ik}\,a_{kj}}{\displaystyle\sum_{\ell \in \CC_i} a_{k\ell}}
    \right)}_{\text{numerator for } w_{ij}}\;e_j \;\approx\; 0.
\]
Solving for $e_i \approx \sum_{j \in \CC_i} w_{ij}\,e_j$ gives the
interpolation weights:
\begin{equation}
  \label{eq:weights}
  \boxed{
    w_{ij} = -\frac{1}{\hat{a}_{ii}}
    \left(a_{ij} + \sum_{k \in D_i^s}
      \frac{a_{ik}\,a_{kj}}{\displaystyle\sum_{\ell \in \CC_i} a_{k\ell}}
    \right),
    \qquad j \in \CC_i.
  }
\end{equation}

\subsection{Properties}

\begin{enumerate}
  \item \textbf{Constant preservation.}  If row~$i$ of~$A$ has zero row
    sum ($\sum_j a_{ij} = 0$), then $\sum_{j \in \CC_i} w_{ij} = 1$.
    This holds for interior nodes of a Laplacian and ensures that the
    constant vector lies in the range of~$P$.

    \begin{proof}
      When $\sum_j a_{ij} = 0$, we have
      $a_{ii} = -\sum_{j \ne i} a_{ij}$, and consequently
      $\hat{a}_{ii} = -\sum_{j \in \CC_i} a_{ij} - \sum_{k \in D_i^s} a_{ik}$.
      Similarly, $\sum_{\ell \in \CC_i} a_{k\ell} = -a_{kk} - \sum_{m \notin \CC_i} a_{km}$
      for each $k$.  Summing~\eqref{eq:weights} over $j \in \CC_i$ and
      using $\sum_{j \in \CC_i} a_{kj} = \sum_{\ell \in \CC_i} a_{k\ell}$
      (the denominator) gives
      $\sum_j w_{ij} = -(\sum_{j \in \CC_i} a_{ij} + \sum_{k \in D_i^s} a_{ik}) / \hat{a}_{ii} = 1$.
    \end{proof}

  \item \textbf{Non-negative weights.}  For M-matrices, $a_{ij} \le 0$
    for $j \ne i$ and $a_{ii} > 0$, so $\hat{a}_{ii} > 0$, the
    numerator in~\eqref{eq:weights} is non-positive, and therefore
    $w_{ij} \ge 0$.

  \item \textbf{Sparsity.}  Each F-point row of~$P$ has at most
    $|S_i \cap \CC|$ entries (the number of strong C-neighbors).
    The total nonzeros in~$P$ equal $|\CC| + \sum_{i \in \FF} |\CC_i|$.
\end{enumerate}

\subsection{Complexity}

For each F-point~$i$, the dominant cost is the inner loop over row~$k$
of~$A$ for each $k \in D_i^s$, giving $O(|D_i^s| \cdot d_{\max})$
where $d_{\max}$ is the maximum row length of~$A$.  Summed over all
F-points, the total cost is $O(\mathrm{nnz}(A) \cdot d_{\max})$ in the
worst case, but $O(\mathrm{nnz}(A))$ for bounded-stencil matrices.

\subsection{Algorithm Summary}

\begin{algorithm}[H]
\caption{Standard Interpolation for F-point $i$}
\label{alg:interp}
\begin{algorithmic}[1]
  \State Identify $\CC_i$, $D_i^s$, $D_i^w$ from $S_i$ and the C/F splitting
  \State $\hat{a}_{ii} \gets a_{ii} + \sum_{m \in D_i^w} a_{im}$
  \For{each $j \in \CC_i$}
    \State $w_{ij} \gets a_{ij}$ \hfill \Comment{direct connection}
  \EndFor
  \For{each $k \in D_i^s$}
    \State $d \gets \sum_{\ell \in \CC_i} a_{k\ell}$
    \If{$|d| < \varepsilon$}
      \State $\hat{a}_{ii} \gets \hat{a}_{ii} + a_{ik}$ \hfill
        \Comment{no common C-points: lump to diagonal}
    \Else
      \For{each $\ell \in \CC_i$ with $a_{k\ell} \ne 0$}
        \State $w_{i\ell} \gets w_{i\ell} + a_{ik}\,a_{k\ell}\,/\,d$
          \hfill \Comment{redistribute proportionally}
      \EndFor
    \EndIf
  \EndFor
  \For{each $j \in \CC_i$}
    \State $w_{ij} \gets -\,w_{ij}\,/\,\hat{a}_{ii}$ \hfill
      \Comment{normalize}
  \EndFor
\end{algorithmic}
\end{algorithm}

%======================================================================
\section{Two-Level Cycle}
%======================================================================

With $P$ constructed, the standard two-level correction scheme is:
\begin{enumerate}
  \item \textbf{Pre-smooth:} Apply $\nu_1$ sweeps of relaxation to
    $Ax = b$.
  \item \textbf{Restrict residual:} $r_c = P^T(b - Ax)$.
  \item \textbf{Coarse solve:} Solve $A_c\,e_c = r_c$ where
    $A_c = P^T A P$.
  \item \textbf{Prolongate and correct:} $x \gets x + P\,e_c$.
  \item \textbf{Post-smooth:} Apply $\nu_2$ sweeps of relaxation.
\end{enumerate}
Applying the coarse solve recursively yields the V-cycle.  The
Galerkin product $A_c = P^T A P$ ensures that the coarse-grid equation
is a variational projection of the fine-grid problem, which is
essential for convergence~\cite{RS1987,TOS2001}.

%======================================================================
\section{Shared Library and Julia Interface}
\label{sec:library}
%======================================================================

The solver is packaged as a shared library (\texttt{lib/libamg.so})
with an opaque-handle C~API suitable for foreign function interfaces.
This allows the AMG+ichol preconditioned CG solver to be called from
any language that supports C~interop (Julia, Python, Fortran, etc.)
without linking individual object files.

\subsection{Library API}

The wrapper (\texttt{src/amg\_wrapper.c}) exposes five functions:

\begin{center}
\begin{tabular}{lp{7.5cm}}
  \toprule
  Function & Description \\
  \midrule
  \verb|amg_solver_create| &
    Accepts 0-based CSR arrays $(n, \mathrm{nnz}, \mathtt{ia}, \mathtt{ja},
    \mathtt{val})$ plus parameters $(\theta, \mathtt{min\_size},
    \mathtt{max\_levels}, \nu)$.  Copies the input, symmetrizes
    $A \leftarrow (A + A^T)/2$, builds the AMG hierarchy, and computes
    the ichol(0) factor.  Returns an opaque \verb|void*| handle. \\
  \verb|amg_solver_solve| &
    Runs PCG with the AMG+ichol preconditioner.  Takes the handle,
    right-hand side~$b$, solution vector~$x$ (zeroed on entry),
    tolerance and maximum iterations.  Returns the iteration count
    (negative if not converged). \\
  \verb|amg_solver_size| & Returns $n$, the matrix dimension. \\
  \verb|amg_solver_num_levels| & Returns the number of AMG levels. \\
  \verb|amg_solver_free| & Releases all memory associated with the handle. \\
  \bottomrule
\end{tabular}
\end{center}

Internally, the handle points to a struct holding the symmetrized
\texttt{dcsr\_mat}~$A$, the \texttt{amg\_data} hierarchy, and the
\texttt{ichol\_factor}~$L$.  The preconditioner callback reuses
\texttt{amg\_ichol\_precond} from the core library.

\subsection{Build}

The \texttt{Makefile} compiles all library sources in \texttt{src/}
with \texttt{-fPIC} and links them into a shared library:
\begin{verbatim}
  make            # produces lib/libamg.so and test_amg
  make clean      # removes objects, test_amg, and lib/libamg.so
\end{verbatim}
The test driver \texttt{test\_amg} links against \texttt{lib/libamg.so}
with an \verb|$ORIGIN/lib| rpath, so no \texttt{LD\_LIBRARY\_PATH}
configuration is needed.

\subsection{Julia Interface}

The script \texttt{solve\_amg.jl} calls \texttt{libamg.so} via Julia's
built-in \texttt{ccall} mechanism, using only standard library modules
(\texttt{SparseArrays}, \texttt{LinearAlgebra}, \texttt{Random},
\texttt{Printf}).

\paragraph{Usage.}
\begin{verbatim}
  julia solve_amg.jl <matrix_file> [threshold=0.25] [nu=1]
\end{verbatim}

\paragraph{CSR conversion.}
Julia stores sparse matrices in compressed sparse column (CSC) format.
To obtain 0-based CSR arrays for the C~library, the script transposes
the matrix: $\texttt{SparseMatrixCSC}(A^T)$ produces CSC of~$A^T$,
which is identical to CSR of~$A$.  Subtracting~1 from
\texttt{colptr} and \texttt{rowval} converts from Julia's 1-based
to the library's 0-based indexing.

\paragraph{Library resolution.}
Julia's \texttt{ccall} requires the library path to be a compile-time
constant.  The script declares
\begin{verbatim}
  const LIBAMG = joinpath(@__DIR__, "lib", "libamg.so")
\end{verbatim}
so that \texttt{@\_\_DIR\_\_} resolves to the directory containing the
script (the project root), and the library is found at
\texttt{lib/libamg.so} without any environment variable configuration.

\paragraph{Data flow.}
\begin{enumerate}
  \item Read the COO file into a Julia sparse matrix and symmetrize.
  \item Convert to 0-based CSR and call \texttt{amg\_solver\_create}.
  \item Generate a random right-hand side $b \in \R^n$.
  \item Call \texttt{amg\_solver\_solve} to fill the solution~$x$.
  \item Compute the relative residual $\|b - Ax\| / \|b\|$ in Julia
    as an independent verification.
  \item Call \texttt{amg\_solver\_free} to release C-side memory.
\end{enumerate}

%======================================================================
\bibliographystyle{plain}
\begin{thebibliography}{5}

\bibitem{RS1987}
  J.~W.~Ruge and K.~St\"uben,
  ``Algebraic multigrid,''
  in \emph{Multigrid Methods} (S.~F.~McCormick, ed.),
  Frontiers in Applied Mathematics, vol.~3,
  SIAM, Philadelphia, 1987, pp.~73--130.

\bibitem{Stuben2001}
  K.~St\"uben,
  ``A review of algebraic multigrid,''
  \emph{Journal of Computational and Applied Mathematics},
  vol.~128, no.~1--2, pp.~281--309, 2001.

\bibitem{TOS2001}
  U.~Trottenberg, C.~W.~Oosterlee, and A.~Sch\"uller,
  \emph{Multigrid},
  Academic Press, 2001.

\bibitem{Saad2003}
  Y.~Saad,
  \emph{Iterative Methods for Sparse Linear Systems},
  2nd~ed., SIAM, 2003.

\bibitem{HensonYang2002}
  V.~E.~Henson and U.~M.~Yang,
  ``{BoomerAMG}: A parallel algebraic multigrid solver and
  preconditioner,''
  \emph{Applied Numerical Mathematics},
  vol.~41, no.~1, pp.~155--177, 2002.

\end{thebibliography}

\end{document}
